
The statements \ref{cdp:st:C.2.CR.1}--\ref{cdp:st:C.2.CR.6} below retain the
proved local and global reductions. The later
\ref{cdp:st:C.2.UN.7} proves that the contracted irrational
configuration has $e(S)\ge2$ and $\pi_1(S)=1$.
Together with the CDP transgression, or the zero
Euler conclusion \ref{cdp:st:C.2.MH.3}, this excludes the
configuration entirely under the CDP hypotheses.
The unrestricted genus assertion is separately
proved in \ref{cdp:st:C.2.UN.10}.

\begin{proposition}[Topology and coefficients in the contracted-section case]\label{cdp:st:C.2.CR.1}
Let $Y$ be a smooth compact complex threefold with
$b_1(Y)=b_2(Y)=0$, and let $f:Y\to\PP^1$ be a proper
holomorphic surjection with connected fibres and smooth
general fibre a complex two-torus. Let $mS$
be a full fibre with irreducible nonnormal reduced support, and
write $h=\eta\pi:Z\to T\to S$ for its normalization and minimal
normalization resolution. Suppose that $p:Z\to B$ is an
irrational ruling, that the effective divisor in \ref{cdp:st:A.5.3} is
$D=2C+V$, where $C$ is a section and $V$ is vertical, and
that $\pi$ contracts $C$. Put $g=g(B)$, let $\Gamma$ be the
conductor-image scheme, and let $n$ be the number of irreducible
components of $\Gamma_{\rm red}$.

Then the following are exact consequences of these hypotheses:
\[
\begin{gathered}
 H^1(T;\mathbb R)=H^1(T,\OO_T)=H^2(T,\OO_T)=0,
 \qquad \operatorname{length}R^1\pi_*\OO_Z=g,\\
 \pi_1(S)\simeq\pi_1(\Gamma_{\rm red})\simeq\mathbb Z,
 \qquad e(\Gamma_{\rm red})=n,\\
 \chi(\OO_\Gamma)=-1,\qquad
 \chi(\OO_{\Gamma_{\rm red}})\in\{-1,0\},\qquad
 \sum_{x\in\Gamma_{\rm red}}\bigl(\delta_x-r_x+1\bigr)\le1.
\end{gathered}
\tag{C.CR.1}\label{cdp:eq:C.CR.1}
\]
Here $r_x$ is the number of normalization branches at $x$.
The full inverse image $D_{\rm red}$ is a smooth genus-$g$
section with trees of smooth rational curves attached; all
its intersections are ordinary double points. For adjacent
components of coefficients $a,b$ one has
\[
                         \gcd(a-1,b-1)=1.          \tag{C.CR.2}\label{cdp:eq:C.CR.2}
\]
Choose a specified sequence of ruling blowdowns to a bundle
$Z_0\to B$. Write its pushed-forward divisor as
$2C_0+\sum_{j=1}^{v}a_jF_j$, with distinct fibres $F_j$
and $a_j>0$, and let $N_0$ count the reverse blowups whose
centre has multiplicity one in the current divisor. Then
\[
                     e(S)=2-2g-v+N_0+n.           \tag{C.CR.3}\label{cdp:eq:C.CR.3}
\]
These necessary equations are retained in the Euler
calculation. \ref{cdp:st:C.2.UN.7} later excludes their simultaneous
geometric occurrence under the stated CDP hypotheses.
\end{proposition}
\begin{proof}
The finite-duality calculation (A.RT.5), whose proof does
not assume rationality, gives $H^2(T,\OO_T)=0$. The
degree-one Leray sequence is consequently
\[
 0\longrightarrow H^1(T,\OO_T)\longrightarrow H^1(Z,\OO_Z)
 \longrightarrow H^0(T,R^1\pi_*\OO_Z)\longrightarrow0.
\tag{C.CR.4}\label{cdp:eq:C.CR.4}
\]
The middle vector space has dimension $g$. At the point
to which $C$ is contracted, formal functions and the
successive infinitesimal exceptional fibres give a
surjection from the completed stalk of $R^1\pi_*\OO_Z$
onto the $H^1$ of the scheme-theoretic exceptional fibre.
Indeed the restriction maps on $H^1$ of these thickenings
are surjective, since their kernels are supported in
dimension at most one. The completed stalk equals the
stalk, because this higher direct image has finite support.
Restriction from the exceptional fibre to $C$ is likewise
surjective on $H^1$. Thus its length is at least
$h^1(C,\OO_C)=g$, whereas \eqref{cdp:eq:C.CR.4} bounds the total length
by $g$. Equality proves the coherent assertions in the
first line of \eqref{cdp:eq:C.CR.1}.

Connectedness of resolution fibres makes the constant-sheaf
Leray map $H^1(T;\mathbb R)\to H^1(Z;\mathbb R)$ injective.
Restriction from $Z$ to its section $C$ is an isomorphism
on $H^1$, because a ruled bundle has simply connected
fibres and point blowups do not change $H^1$. The image of
$H^1(T;\mathbb R)$ restricts to zero on the contracted
curve $C$. This proves its vanishing.

We describe the actual divisor comparisons. At every ruling
blowdown $\sigma:W\to W'$ with exceptional curve $E$, let
$a$ be the coefficient of $E$ and let $d$ be the multiplicity
of the pushed-forward boundary at the centre. The identity
\[
 K_W+D_W=\sigma^*(K_{W'}+D_{W'})+(1+a-d)E
\tag{C.CR.5}\label{cdp:eq:C.CR.5}
\]
and numerical triviality give $a=d-1$. The minimal model
boundary is therefore $2C_0+\sum a_jF_j$. Its reduced
support is the section with disjoint rational fibre leaves.
Every reverse blowup preserves the asserted simple-normal-
crossings tree structure: a positive exceptional coefficient
adds a leaf at a smooth boundary point or subdivides an edge
at a boundary node; a zero exceptional coefficient leaves
the reduced boundary isomorphic to its predecessor. A centre
off the current boundary would create coefficient $-1$,
which persists on its strict transform and contradicts the
effectiveness of the original divisor. In particular
$D_{\rm red}$ is connected, and its fundamental group maps
isomorphically to that of $Z$ by inclusion, with both
identified with $\pi_1(B)$ through the section.

This also proves \eqref{cdp:eq:C.CR.2} without discarding coefficients.
Initially every edge has coefficients $(2,a_j)$ and hence
coprime shifted coefficients. A node blowup inserts the
coefficient $a+b-1$, so its shifted coefficient is
$(a-1)+(b-1)$; the two new edge gcds equal the old gcd.
A positive-coefficient blowup at a smooth point of a
coefficient-$a$ boundary creates coefficient $a-1$ and
an edge with shifted coefficients $a-1,a-2$, whose gcd
is one. These are all the possible operations.

On a minimal normalization resolution, a coefficient-zero
exceptional curve $E$ has $K_Z\cdot E\ge0$. Since
$(K_Z+D)\cdot E=0$, it follows that $D\cdot E=0$.
Such a curve cannot meet $D_{\rm red}$. Connectedness of
exceptional fibres and of $D_{\rm red}$ now gives the same
support equality as in A.5.16:
\[
 D_{\rm red}=h^{-1}(\Gamma_{\rm red}).
\tag{C.CR.6}\label{cdp:eq:C.CR.6}
\]
Any exceptional fibre outside this set would correspond to
an isolated normal hypersurface singularity of $S$, because
the normal conductor ideal is reflexive and has no isolated
support. \ref{cdp:st:A.5.17} excludes every such point in the actual
torus smoothing. Thus $h$ is an isomorphism off the two
closed sets in \eqref{cdp:eq:C.CR.6}.

It follows that $S$ is the topological pushout of
$D_{\rm red}\hookrightarrow Z$ and
$D_{\rm red}\to\Gamma_{\rm red}$. This is the actual
pushout: its map to $S$ is a continuous bijection from a
compact space to a Hausdorff space. Triangulation of the
pair makes the first map a cofibration. Van Kampen applied
to its double mapping cylinder therefore gives
\[
 \pi_1(S)=\pi_1(Z)*_{\pi_1(D_{\rm red})}
                         \pi_1(\Gamma_{\rm red})
          \simeq\pi_1(\Gamma_{\rm red}).
\tag{C.CR.7}\label{cdp:eq:C.CR.7}
\]
Indeed the first amalgamating map is an isomorphism, and
the second map is trivial, since the section $C$ generates
$\pi_1(D_{\rm red})$ and $h(C)$ is a point. Every component
of $\Gamma_{\rm red}$ is dominated by a vertical rational
curve, so its normalization is $\PP^1$. Its fundamental
group is the free group of rank $b_1(\Gamma_{\rm red})$.

For the exact local comparison with the nearby torus, take
a triangulated regular neighbourhood retracting onto $S$
and put a sufficiently small full fibre tube inside it.
Perturbing loops off the divisor gives a surjection from
the fundamental group of its punctured torus bundle onto
that of $S$. The former is $\mathbb Z^4\rtimes\mathbb Z$,
which is metabelian. A free group of rank at least two
is not metabelian, since it surjects onto $S_4$.
Consequently the rank in \eqref{cdp:eq:C.CR.7} is zero or one. The
nonzero integral class on every support supplied by \ref{cdp:st:D.1.9}
excludes zero. This proves the fundamental-group assertion.

Normalization of the rational image components gives the
exact identities
\[
 b_1(\Gamma_{\rm red})
   =\sum_x(r_x-1)-n+1=1,\qquad
 e(\Gamma_{\rm red})=2n-\sum_x(r_x-1)=n.
\tag{C.CR.8}\label{cdp:eq:C.CR.8}
\]
Put $\epsilon=h^0(S,\omega_S)\in\{0,1\}$.
The vanishings already proved, the normalization exact
sequence and conductor duality give
\[
 h^0(\eta_*\OO_T/\OO_S)=1+\epsilon,\qquad
 h^1(\eta_*\OO_T/\OO_S)=\epsilon,
 \qquad\chi(\OO_\Gamma)=-1.
\tag{C.CR.9}\label{cdp:eq:C.CR.9}
\]
The nilradical sequence and curve duality give
$\chi(\OO_{\Gamma_{\rm red}})\ge-1$ exactly as in
C.2.VII.2: for $\epsilon=0$ the relevant twisted $H^1$
has dimension at most one, and for $\epsilon=1$ the twist
is trivial and connectedness supplies $h^0=1$.
On a connected curve with rational normalizations,
$\chi(\OO_{\Gamma_{\rm red}})=n-\sum_x\delta_x$.
Together with \eqref{cdp:eq:C.CR.8} and $\delta_x\ge r_x-1$, this
proves the remaining assertions of \eqref{cdp:eq:C.CR.1}.

Finally $e(Z_0)=4-4g$ and
$e((D_0)_{\rm red})=2-2g+v$. A reverse blowup of
multiplicity at least two increases the Euler numbers
of both the surface and the reduced boundary by one,
leaving their difference unchanged. A multiplicity-one
centre leaves the boundary isomorphic to its predecessor
and increases the complement Euler number by one.
Therefore
\[
 e(Z\setminus D_{\rm red})=2-2g-v+N_0.
\]
Equation \eqref{cdp:eq:C.CR.6}, the excluded off-conductor exceptionals,
and additivity of Euler number now give \eqref{cdp:eq:C.CR.3}.
\end{proof}


